% CDS547 Group Project Progress Report
% Guideline: 5--8 pages excluding appendices and references.
% If your instructor requires the official Overleaf template, copy these sections into that project.
\documentclass[11pt,a4paper]{article}

\usepackage[margin=1in]{geometry}
\usepackage{graphicx}
\usepackage{booktabs}
\usepackage{hyperref}
\usepackage{enumitem}
\usepackage{amsmath}
\usepackage{cite}

\title{\textbf{Group Project Progress Report}\\
\large CDS547 -- Introduction to Large Language Models\\
\large Retrieval-Augmented Generation for Domain-Specific Document QA}
\author{Group Name / Number \quad [Replace with member names]\\
\textit{[Institution]}}
\date{March 2026}

\begin{document}
\maketitle

\begin{abstract}
This progress report documents the midpoint status of our CDS547 group project on \emph{retrieval-augmented generation (RAG)} for domain-specific document question answering. We summarize completed work in literature review, system design, and preliminary implementation; present early observations and partial evaluation; discuss challenges and deviations from the original plan; and outline responsibilities and next steps toward the final deliverables.
\end{abstract}

\section{Introduction and Project Recap}
Our proposal (submitted February~2026) committed to building a RAG pipeline over a curated document corpus and comparing it to direct large language model (LLM) generation without retrieval. The motivation is to reduce factual hallucinations and improve traceability by conditioning generation on retrieved passages \cite{lewis2020retrieval}.

\textbf{Original objectives} included: (1)~implementing ingestion, chunking, embedding, and retrieval; (2)~integrating an LLM for answer generation with optional citations; (3)~evaluating RAG vs.\ no-retrieval baselines; and (4)~analyzing retrieval design choices (chunk size, top-$k$, retrieval modality). This report assesses progress against these goals at the midpoint deadline (March~11).

\section{Technical Progress}

\subsection{Research and literature review}
Completed survey reading on transformer-based encoders \cite{vaswani2017attention,devlin2019bert}, few-shot behavior of large LMs \cite{brown2020language}, and the RAG paradigm \cite{lewis2020retrieval}. The group aligned on definitions of \emph{faithfulness} and \emph{grounding}: answers should be supported by retrieved text, and unsupported claims should be avoided or explicitly hedged. We identified evaluation approaches: reference-based metrics where gold answers exist, and rubric-based or model-assisted grading for open-ended QA.

\subsection{Implementation}
\textbf{Data and preprocessing.} We selected [REPLACE: e.g., a subset of course PDFs / a public dataset such as Natural Questions or a domain corpus]. Documents were converted to plain text, cleaned, and segmented into chunks of [REPLACE: e.g., 256--512 tokens] with [REPLACE: fixed / sliding] overlap to preserve context across boundaries.

\textbf{Retrieval.} We implemented [REPLACE: dense vector store + embedding model / hybrid sparse+dense]. The index supports top-$k$ retrieval ($k \in \{3,5,10\}$ for upcoming ablations). Initial sanity checks (keyword search and nearest-neighbor spot checks) suggest chunks are on-topic for pilot queries.

\textbf{Generation.} The pipeline formats retrieved chunks into a prompt template instructing the model to answer only from context and to cite passage indices when possible. We support [REPLACE: OpenAI API / open-weight model via Hugging Face / other].

\textbf{Baseline.} A \emph{no-retrieval} prompt with the same question and instructions but empty context is implemented for controlled comparison.

\subsection{Experimentation (preliminary)}
We drafted an evaluation set of [REPLACE: N] questions with [REPLACE: short answers or rubric criteria]. Early runs indicate:
\begin{itemize}[noitemsep]
    \item RAG answers often include phrasing aligned with source passages, improving perceived faithfulness on factual queries.
    \item Failure modes observed: irrelevant retrieval when queries are ambiguous; ``correct-sounding'' answers when retrieval misses the decisive sentence; sensitivity to chunk boundaries.
    \item Latency and API cost scale with context length; we are logging tokens per query for the final cost--quality analysis.
\end{itemize}
Quantitative tables are in preparation once the evaluation set is frozen and all baselines are run with fixed seeds and API versions.

\section{Critical Analysis and Early Results}
\textbf{Strengths.} RAG provides an auditable path from claim to evidence, which supports debugging and aligns with course themes on responsible deployment. Pilot qualitative review suggests fewer unsupported statements than the no-retrieval baseline on our [REPLACE: domain] questions.

\textbf{Limitations.} (1)~Retrieval quality caps generation quality; errors propagate. (2)~Small corpora may yield repetitive or incomplete coverage. (3)~LLM-assisted evaluation can bias scores; we plan to disclose this and include human spot checks on a subset.

\textbf{Plan adjustments.} We originally considered [REPLACE: optional fine-tuning]. Given time and compute constraints, we narrowed scope to \emph{prompting + RAG} without fine-tuning, and added an explicit milestone to freeze the evaluation set by week~[REPLACE] to avoid leakage from iterative debugging.

\section{Project Management and Timeline}
\textbf{Roles} follow the proposal: team lead coordinates milestones; researcher maintains bibliography and related work; developer owns repository structure and reproducibility (README, environment file); evaluator owns metrics spreadsheets and result plots.

\textbf{Schedule check.} Literature review and core pipeline are largely on track. Experimentation is slightly behind due to [REPLACE: e.g., data access / API quota / chunking bugs]; mitigation includes [REPLACE: e.g., smaller pilot corpus first, shared API key rotation, pair programming on indexer].

Table~\ref{tab:timeline} summarizes revised near-term tasks.

\begin{table}[h]
\centering
\caption{Near-term tasks (illustrative --- customize dates).}
\label{tab:timeline}
\begin{tabular}{@{}ll@{}}
\toprule
Target week & Task \\
\midrule
Current & Freeze chunking + embedding choices; lock eval questions v1 \\
+1 week & Full baseline runs (RAG vs.\ no retrieval); log all outputs \\
+2 weeks & Ablation on top-$k$ and chunk size (subset) \\
+3 weeks & Draft results section; begin presentation storyboard \\
\bottomrule
\end{tabular}
\end{table}

\section{Next Steps and Risks}
\textbf{Next steps.} Complete end-to-end runs with frozen configuration; compute quantitative metrics; add error taxonomy (retrieval vs.\ generation failures); prepare figures for final report; rehearse demo with one held-out document set.

\textbf{Risks.} API instability or pricing changes; inconsistent document formatting hurting chunk quality; evaluation set too small for strong claims. \textbf{Contingency:} emphasize qualitative case studies and clear ablation on one axis (e.g., top-$k$ only) if time is constrained.

\section*{Acknowledgments}
[Optional: thank TA/instructor or cite open datasets/tools.]

\bibliographystyle{plain}
\begin{thebibliography}{9}

\bibitem{vaswani2017attention}
A.~Vaswani et al.
\newblock Attention is all you need.
\newblock In \emph{NeurIPS}, 2017.

\bibitem{devlin2019bert}
J.~Devlin et al.
\newblock BERT: Pre-training of deep bidirectional transformers for language understanding.
\newblock In \emph{ACL}, 2019.

\bibitem{brown2020language}
T.~Brown et al.
\newblock Language models are few-shot learners.
\newblock In \emph{NeurIPS}, 2020.

\bibitem{lewis2020retrieval}
P.~Lewis et al.
\newblock Retrieval-augmented generation for knowledge-intensive NLP tasks.
\newblock In \emph{NeurIPS}, 2020.

\end{thebibliography}

\appendix
\section*{Appendix A: Reproducibility Checklist (optional)}
[Repository link (private if needed), Python version, key package versions, embedding model ID, LLM model name and date, chunking parameters.]

\end{document}
